% !TEX root = ../nc-quadric-classification.tex
% 2026-11-21 / v1.39
\section{noncubic 境界の適合根と Hochschild 作用}
\label{sec:noncubic-hh-roots}


\paragraph{結論と範囲.}
中央二重平方表示の二つの根対を交換する\emph{trace 零の}線型対合を選ぶと、
対応する適合 Serre 根は有限 foundation の $\operatorname{HH}^1,\operatorname{HH}^2$ に恒等に作用する。
ただし根対が一致するときは対角の交換対合を選ぶ。
trace 零だけでは十分でないことと、選択条件の具体的な行列式を後で示す。
この計算は noncubic 表示の一次元の変形だけを扱うものではない。
foundation の全変形接空間、すなわち根対が異なる場合の三次元、
一致する場合の四次元を扱う。
従って既存の Artin 根持上げ補題を適用すると、指定した良い根は
全ての Artin Abel 変形へ持ち上がり、その Morita 類は一意である。
さらに有限 foundation の各 noncubic 点には、
この根と実際の Serre 平方同型を持つ有限型 étale 近傍が存在する。
同じ根に対する AS4 被覆の横断変形への持上げと、有限型基底上の全相対認識は
本節の結論に含めない。
本節の計算の新規性は未確認である。

\subsection{四項 foundation の二項移動}

$k$ を標数零の代数閉体とする。
full geometric helix の四項を $E_0,E_1,E_2,E_3$ とし、
\[
A=\Hom(E_0,E_1),\quad B=\Hom(E_1,E_2),\quad
C=\Hom(E_2,E_3),\quad D=\Hom(E_0,E_2)
\]
と置く。$A,B,C$ は二次元、$D$ は四次元である。
$W=\Hom(E_0,E_3)$ は六次元であり、
\[
F:A\otimes B\longrightarrow D,\qquad
G:C\otimes D\longrightarrow W
\]
が foundation の積を与える。
因子の並べ替えを省略して書く。
$G$ は全射とし、$K=\ker G$ は二次元とする。
包含 $K\subset C\otimes D$ を curry した写像を
\[
j:D^{\vee}\longrightarrow C\otimes K^{\vee}
\]
と書く。さらに
\begin{equation}\label{eq:nchh-delta}
\delta=(-F^{\vee},j):D^{\vee}\longrightarrow
H=(A^{\vee}\otimes B^{\vee})\oplus(C\otimes K^{\vee})
\end{equation}
は単射とする。これらは、以下で使う noncubic 点の近傍で開条件である。
$\mathscr L=\operatorname{im}\delta\in\Gr(4,H)$ と置く。

\begin{lemma}[二項移動の有限線型代数]\label{lem:nchh-gale}
$E_{i+4}=S^{-1}E_i[2]$ により延長し、四項窓を
$(E_2,E_3,E_4,E_5)$ へ移す。
対応するデータは
\begin{align}
A'&=C,& B'&=K^{\vee},& C'&=A,\label{eq:nchh-new-small}\\
D'&=H/\mathscr L,&
F'&:C\otimes K^{\vee}\longrightarrow D',\label{eq:nchh-new-D}\\
K'&=\operatorname{im}\bigl(B^{\vee}\longrightarrow A\otimes D'\bigr),\label{eq:nchh-new-K}
\end{align}
である。\eqref{eq:nchh-new-K} は $A$ の coevaluation と第一 summand の商写像の合成である。
$K'\simeq B^{\vee}$ となる開近傍で、この新しい四平面は
\begin{equation}\label{eq:nchh-gale}
\mathscr L'=\{(-Y,X):(X,Y)\in\mathscr L^{\perp}\}
 \subset(C^{\vee}\otimes K)\oplus(A\otimes B)
\end{equation}
である。ここで $\perp$ は $H$ と $H^{\vee}$ の標準評価による annihilator を表す。
この記述は任意の局所 Artin 変形にも同じ式で適用できる。
\end{lemma}

\begin{proof}
全証明は相対モジュライ論文 \cite[\texttt{noncubic_hh_roots}]{QuadricModuli2026} に置く。
\end{proof}

\begin{remark}[全変形空間を使うこと]
$G$ の商の基底を忘れ、代わりに $K$ の基底を選ぶと、
foundation のデータは $(F,j)$ である。
$\delta$ が単射である開部分で $D$ の基底変更を除くことは、
四平面 $\mathscr L\subset H$ を取ることに等しい。
従って全 foundation 変形の局所模型は
$\Gr(4,8)$ の開部分を
\[
\widetilde G=\GL(A)\times\GL(B)\times\GL(C)\times\GL(K)
\]
で割ったものとなる。
第一 summand への作用は $g_A^{-t}\otimes g_B^{-t}$、
第二 summand への作用は $g_C\otimes g_K^{-t}$ である。
自明に作用する三次元中心を除いた群を $G_{\mathrm{eff}}$ と書く。
その次元は十三で、Grassmannian の次元は十六である。
この模型には noncubic 表示を保たない全ての横断方向が含まれている。
\end{remark}

\subsection{中央二重平方と良い根の選択}

\begin{lemma}[trace 零の交換対合]\label{lem:nchh-exchange}
二つの非退化二元二次形式 $f_1,f_2$ に対し、座標とそれぞれの平方根変数の
定数倍を選び、
\begin{equation}\label{eq:nchh-normal-forms}
f_1=uv,\qquad f_2=\ell_1\ell_2,
\qquad
\begin{pmatrix}\ell_1\\\ell_2\end{pmatrix}
=L_0\begin{pmatrix}u\\v\end{pmatrix},\qquad
L_0=\begin{pmatrix}\alpha&\beta\\\gamma&-\alpha\end{pmatrix},\quad
\alpha^2+\beta\gamma=1
\end{equation}
とできる。二根対が一致する場合にも $L_0=\operatorname{diag}(1,-1)$ と選べる。
\end{lemma}
\begin{proof}
全証明は相対モジュライ論文 \cite[\texttt{noncubic_hh_roots}]{QuadricModuli2026} に置く。
\end{proof}

以下、
\[
\mathcal B=k[u,v]\langle c,d\rangle/(c^2-uv,d^2-\ell_1\ell_2)
\]
と置く。$u,v$ は中央である。
\begin{equation}\label{eq:nchh-sigma}
\sigma(u,v)=L_0(u,v),\qquad\sigma(c)=d,\qquad\sigma(d)=-c
\end{equation}
は次数付き自己同型で、$\sigma^2=\iota$、ただし
$\iota(u)=u,\iota(v)=v,\iota(c)=-c,\iota(d)=-d$ である。
命題~\ref{prop:noncubic-relative-serre-on-perfect} の双加群計算から
$S=(-)_\iota(-2)[2]$ であり、
\begin{equation}\label{eq:nchh-good-root}
s=(-)_\sigma(1),\qquad s^2=S^{-1}[2]
\end{equation}
は適合根である。同じ根の AS4 表示は定理~\ref{thm:ore-anticommuting-central} の Ore/Zhang 構成が与える。

\begin{proposition}[根の作用と交代 Gale 双対]\label{prop:nchh-root-gale}
\eqref{eq:nchh-good-root} の Hochschild 作用は、前節の Grassmann 模型で
\begin{equation}\label{eq:nchh-involution}
\tau(\mathscr L)=J\mathscr L^{\perp},\qquad
J(X,Y)=(D_0Y,-D_0X),\qquad D_0=\operatorname{diag}(1,-1)
\end{equation}
の微分により与えられる。$X,Y$ を二次正方行列として書いている。
$J:H^{\vee}\to H$ は交代同型であり、中央二重平方の点 $\mathscr L$ は
$\tau$ の固定点、すなわち対応する交代形式の Lagrangian 四平面である。
\end{proposition}

\begin{proof}
全証明は相対モジュライ論文 \cite[\texttt{noncubic_hh_roots}]{QuadricModuli2026} に置く。
\end{proof}

\subsection{全 Hochschild 接空間への作用}

\begin{theorem}\label{thm:noncubic-hh-trivial-root}
任意の noncubic $\mathcal H_2$ 対象の有限 foundation を $\Gamma$ とする。
補題~\ref{lem:nchh-exchange} の trace 零の交換対合を用い、
根対が一致する場合には $L_0=D_0$ と選ぶ。
この選択による適合根 $s$ は
\[
s_*|_{\operatorname{HH}^1(\Gamma)}=1,\qquad s_*|_{\operatorname{HH}^2(\Gamma)}=1
\]
を満たす。
二つの根対が異なる場合は $\dim\operatorname{HH}^1=0$, $\dim\operatorname{HH}^2=3$、
一致する場合はそれぞれ $1,4$ である。
\end{theorem}

\begin{proof}
全証明は相対モジュライ論文 \cite[\texttt{noncubic_hh_roots}]{QuadricModuli2026} に置く。
\end{proof}

\begin{corollary}[全 Artin 変形への根の持上げ]\label{cor:noncubic-root-all-artin}
前定理の根 $s$ は、閉ファイバーを固定した任意の局所 Artin Abel 変形へ、
相対 Serre 平方根として持ち上がる。
その Morita 自己同値の同型類は一意である。
この一意性は平方同型自体の一意な選択を意味しない。
\end{corollary}
\begin{proof}
全証明は相対モジュライ論文 \cite[\texttt{noncubic_hh_roots}]{QuadricModuli2026} に置く。
\end{proof}

\subsection{全 foundation 変形の Lagrangian 正規形}
\label{subsec:noncubic-lagrangian-slice}

前の接計算は、有限型の具体的な局所表示に強められる。
この表示は foundation の全横断方向を含むが、
それだけで相対 $\mathcal H_2$ の認識を与えるものではない。

四つの二次元空間をそれぞれ $k^2$ と同定し、
第\ref{sec:noncubic-hh-roots}節の $H$ と $J$ を使う。
$\omega$ を行列 $J^{-1}$ の非退化交代形式とする。
$X=\Gr(4,H)$、$L=\operatorname{LG}(4,H,\omega)$ と置けば、
$L$ は $\tau(\mathscr L)=J\mathscr L^\perp$ の固定部分スキームである。
基底変更群の作用は
\[
 \rho(g_A,g_B,g_C,g_K)
 =\operatorname{diag}(g_A^{-t}\otimes g_B^{-t},
                      g_C\otimes g_K^{-t})
\]
である。$X$ 上で自明に作用する中心を
\[
 C_0=\{(aI,bI,cI,dI):a^{-1}b^{-1}=cd^{-1}\}
       \simeq\mathbf G_m^3,
 \qquad G=\widetilde G/C_0
\]
と書く。

\begin{proposition}[Lagrangian な有限型 \texorpdfstring{\'etale}{etale} 表示]
\label{prop:noncubic-lagrangian-etale-slice}
群対合
\begin{equation}\label{eq:noncubic-group-involution}
 \widetilde\theta(g_A,g_B,g_C,g_K)
 =(D_0g_CD_0,g_K^{-t},D_0g_AD_0,g_B^{-t})
\end{equation}
は $G$ の対合 $\theta$ を誘導し、
$\tau(g\mathscr L)=\theta(g)\tau(\mathscr L)$ を満たす。
その固定群の単位成分を $G^+=(G^\theta)^\circ$ と書くと、
\[
 G^+\simeq
 (\GL_2\times\GL_2)/(\mathbf G_m\times\mu_2),
 \qquad \dim G^+=7.
\]
ここで中央部分群の埋込みは
$(z,\epsilon)\mapsto(zI,\epsilon z^{-1}I)$ である。

$G$ の Lie 代数の負固有空間を $\mathfrak g^-$ とする。
自然な軌道微分から得る束写像
\begin{equation}\label{eq:noncubic-normal-orbit-map}
 q:\mathfrak g^-\otimes\mathcal O_L
   \longrightarrow N_{L/X}
   \simeq\bigwedge^2\mathscr U^\vee
\end{equation}
を考える。$\mathscr U$ は $L$ 上の tautological 四平面束である。
その行列式の非消滅開を $L^{\mathrm{tr}}$ とすると、
\[
 [L^{\mathrm{tr}}/G^+]\longrightarrow[X/G]
 \tag{LG-slice}\label{eq:noncubic-lagrangian-stack-map}
\]
は表現可能な \texorpdfstring{\'etale}{etale} 射である。
双方のスタックの次元は三である。
$\det q$ は、定数一次元因子を除けば
Pl\"ucker 線束の三乗 $\mathcal O_L(3)$ の切断である。

中央二重平方の trace 零正規形
\eqref{eq:nchh-plane} が $L^{\mathrm{tr}}$ に属する必要十分条件は
$\alpha\ne0$ である。
根対が異なればこの条件は自動的に成立し、
根対が一致するときは補題~\ref{lem:nchh-exchange} の通り
$L_0=D_0$ と選べば成立する。
従ってその近傍の全 foundation 変形は、
\texorpdfstring{\'etale}{etale} 局所に Lagrangian 四平面で表示できる。
同じ結論は、補題~\ref{lem:nchh-gale} の階数条件を課した
foundation の開部分へ制限したときにも成立する。
\end{proposition}

\begin{proof}
全証明は相対モジュライ論文 \cite[\texttt{noncubic_hh_roots}]{QuadricModuli2026} に置く。
\end{proof}

\begin{remark}[この表示が与えるもの]
\label{rem:noncubic-lagrangian-scope}
上の射の像は foundation スタックの開部分であり、
全 noncubic 点を含む。その smooth 性は有限代数表示についての主張である。
実際の Morita 平方根との比較には
系~\ref{cor:noncubic-etale-morita-root} を使い、
必要に応じて二つの \texorpdfstring{\'etale}{etale} 近傍の積を取る。
全正軌道 Hom、相対 heart、同じ根の AS4 被覆の存在を、
この Lagrangian 正規形だけから結論してはならない。
特に \eqref{eq:noncubic-lagrangian-stack-map} の標的を
内在的 $\mathfrak H_2$ に置き換えたとは主張していない。
\end{remark}

\begin{remark}[trace 零だけでは良い根を特徴付けない]
\label{rem:noncubic-offdiagonal-root}
根対一致の場合、$L_0=\left(\begin{smallmatrix}0&1\\1&0\end{smallmatrix}\right)$
も trace 零の交換対合であるが、上の \texorpdfstring{\'etale}{etale} 開には属さない。
対応する根は連結慣性
$u\mapsto zu,v\mapsto z^{-1}v$ を逆元へ移すため、
$\operatorname{HH}^1$ に $-1$ で作用する。
\eqref{eq:nchh-tangent-sequence} の正負分解より、この場合
\[
 (\dim\operatorname{HH}^2_+,\dim\operatorname{HH}^2_-)=(3,1).
\]
従ってこの根は全ての一次 Abel 変形に延長するわけではない。
本稿で全 Artin 延長に使う良い根は、根対一致時には
$L_0=D_0$ を選ぶという条件を含む。
\end{remark}

\subsection{Hochschild 二次に自明な根の étale 代数化}
\label{subsec:hh2-etale-morita-root}

次の命題は noncubic 点に限らない。
根の存在する空間から étale な切片を取ることにより、
Hochschild 一次への作用を仮定せずに根を代数化する。

\begin{proposition}[Hochschild 二次に自明な根の étale 代数化]
\label{prop:hh2-etale-morita-root}
$S$ を標数零の代数閉体 $k$ 上の局所有限型 scheme とし、
$\Lambda$ を $S$ 上の有限局所自由代数とする。
次を仮定する。
\begin{enumerate}
\item $\operatorname{Perf}(\Lambda)$ は相対 smooth proper であり、
任意の基底変更後も
$\operatorname{HH}^{<0}(\Lambda)=0$、
$\operatorname{HH}^0(\Lambda)=\OO_S$ である。
\item 相対 inverse Serre $t=S_\Lambda^{-1}[2]$ は
可逆 perfect 双加群で表され、任意の導来基底変更と可換する。
\item $x\in S(k)$ と
$p_0\in\operatorname{DPic}_k(\Lambda_x)$ が与えられ、
$p_0^2=t_x$、
$(p_0)_*|_{\operatorname{HH}^2(\Lambda_x)}=1$ である。
\end{enumerate}
$\mathcal G\to S$ を Morita 自己同値の étale 局所同型類の代数空間、
$\tau=[t]$ とし、
\[
 \mathcal Z=
 \mathcal G\times_{\mathrm{sq},\mathcal G,\tau}S
\]
を平方根空間とする。
このとき $\mathcal Z\to S$ は $p_0$ で smooth である。
さらに有限型 affine な étale 近傍 $(U,u)\to(S,x)$ 上に、
可逆 perfect $\Lambda_U$-双加群 $P$ と実際の同型
\[
 P\otimes_{\Lambda_U}^{\mathbf L}P\simeq t_U
 \tag{N-root}\label{eq:hh2-etale-actual-root}
\]
が存在し、$[P_u]=p_0$ となる。
追加で $(p_0)_*|_{\operatorname{HH}^1(\Lambda_x)}=1$ なら、
$\mathcal Z\to S$ 自体が $p_0$ で étale であり、
同じ閉根を持つ Artin 持上げの Morita 同型類は一意である。
\end{proposition}

\begin{proof}
全証明は相対モジュライ論文 \cite[\texttt{noncubic_hh_roots}]{QuadricModuli2026} に置く。
\end{proof}

\begin{remark}[Hochschild 一次の役割]
\label{rem:hh2-etale-root-tangent}
同じ閉根を持つ根の一次的な差は
$a\in\operatorname{HH}^1(\Lambda_x)$ で
$a+(p_0)_*a=0$ を満たすものである。
従って平方根空間の相対接空間は
$\ker(1+(p_0)_*)$ である。
Hochschild 一次への恒等作用はこの相対接空間を消し、
根の空間自体を étale にするために使う。
étale 近傍上で一つの根を選ぶ存在定理には不要である。
\end{remark}

\begin{corollary}[全 noncubic 点の有限型 étale Morita 根]
\label{cor:noncubic-etale-morita-root}
$S$ を有限 foundation スタックの有限型 affine chart とし、
$x\in S(k)$ を任意の noncubic $\mathcal H_2$ 対象の
normal foundation に対応する点とする。
定理~\ref{thm:noncubic-hh-trivial-root} の良い根 $s_0$ を選ぶ。
このとき $x$ の有限型 affine étale 近傍 $(U,u)\to(S,x)$ 上に、
$s_0$ を閉 fiber とする可逆 perfect 双加群 $P$ と
実際の平方同型 \eqref{eq:hh2-etale-actual-root} が存在する。
平方根空間自体もこの点で étale である。

四つの普遍 projective 対象 $F_0,\ldots,F_3$ を、
閉 fiber で
$(\OO,\pi M_c,\OO(1),\pi M_d(1))$ に対応させる。
近傍をさらに縮めれば、$P$ による自己同値 $\mathsf s$ は
\[
 \mathsf s F_0\simeq F_2,\qquad
 \mathsf s F_1\simeq F_3
 \tag{N-window}\label{eq:noncubic-etale-root-window}
\]
を満たす。これらの同型は任意の導来基底変更に保たれる。
\end{corollary}

\begin{proof}
全証明は相対モジュライ論文 \cite[\texttt{noncubic_hh_roots}]{QuadricModuli2026} に置く。
\end{proof}

\begin{remark}[有限型近傍でまだ結論していないこと]
\label{rem:noncubic-etale-root-scope}
上の $U$ は有限代数と perfect Morita 根の実際の有限型族である。
この構成だけから相対 Abel heart、
全正軌道 Hom の同じ開近傍での次数零集中、
全軌道環の平坦性・Noetherian 性・AZ 豊富性、
同じ根に対する相対 AS4 被覆は従わない。
従って $U\to\mathfrak H_2$ や
noncubic 点を含む smooth Artin 開近傍を、ここでは構成したとは主張しない。
明示的な noncubic 二形式表示族の効果性と、
cubic を横断するこの étale 近傍の効果性も区別する。
\end{remark}

\subsection{根の有限性、計算の検証と残る問題}

根対が異なる場合、命題~\ref{prop:first-orbit-object} により任意の適合根の
$\OO$ 上の像は同じ内在的対象である。
従って二つの根の差は pointed 自己同値を与え、全根の同型類は
有限な pointed 自己同値群の torsor の部分集合に入る。
一般点でのこの群は既存原稿では $D_8\times C_2$（位数十六）であり、
二つの平方根変数の符号群 $\mu_2^2$ だけに限定してはいけない。
本節は全ての根の作用を恒等と主張せず、上記の trace 零の選択による
一つの根を特定して証明した。
他の根に対する作用は、それとの差である pointed 自己同値の作用との合成である。

補足の厳密有理計算
\texttt{computations/verify\_noncubic\_gale\_root.py} は、非交差、一根共有、一致の各配置で
$\mathscr L$、$J$、Grassmann tangent の十六次元行列、
基底変更軌道の全十六生成元を構成する。
浮動小数点や乱数を使わず、
\[
\begin{array}{c|ccc}
 &\text{非交差}&\text{一根共有}&\text{一致}\\\hline
\dim(\text{orbit}^+,\text{orbit}^-)&(7,6)&(7,6)&(6,6)\\
\dim(\operatorname{HH}^2_+,\operatorname{HH}^2_-)&(3,0)&(3,0)&(4,0)
\end{array}
\]
を確認する。一般の配置への証明はこの三例の計算から外挿せず、
定理~\ref{thm:noncubic-hh-trivial-root} の同変接複体と安定化群の議論による。

\paragraph{同じ根の AS4 被覆は別の問題である.}
持ち上げた同じ根の軌道環の $B_1$ は四次元の $D$ であり、
$B_2=\Hom(E_0,E_4)$ と積は同じ有限 foundation の bar 複体から計算できる。
根を適合するように twist する前の中央二重平方環では、二次双対に正規な二次元の
pencil が存在するが、固定した適合根の $\mathcal B^\sigma$ でも
同じ pencil が残るとは限らない。ambient の関係線の $\sigma$ 不変性が必要である。
非交差・一根共有・一致の三つの有理配置では、三次平坦性を保つ
関係空間の接次元 $18,18,19$ の全方向に対し、選んだ projective 正規線と
その normalizer が一次に一意に持ち上がることを
\texttt{computations/verify\_noncubic\_cover\_tangent.py} の有限線型代数で確認した。
これはこの特定の被覆の一次計算であり、それ自体から高次存在は従わない。
次節では別の中央 AS4 被覆を全配置で構成し、
系~\ref{cor:noncubic-all-artin-h2} で同じ良い根に対する全 Artin 延長を証明する。
有限型近傍の一様な効果性はなお別問題である。
